\subsection{Hasse duality and the Eisenstein bridge}

Let $p\geq5$, put $P=p(p-1)$, and work first on a compactified Katz
modular curve $X$ with auxiliary full level $3$. Let $\omega$
be the Hodge bundle, let $C$ be the cusp divisor, and let
$A\in H^0(X,\omega^{p-1})$ be the Hasse invariant. Its zero divisor
$S$ is reduced and disjoint from $C$. Put $Z=pS$, so that $Z$ is the
Cartier divisor cut out by $A^p$. For every integer $r$ define
\[
 \mathcal Q_r:=\omega^r/A^p\omega^{r-P}
              \simeq \omega^r\vert_Z,
 \qquad Q_r:=H^0(Z,\mathcal Q_r).
\]
The constructions below commute with the tame deck group. Its order is
prime to $p$, so the Reynolds operator is exact.  The canonical pairing
below is deck-invariant, and its restriction to invariant subspaces is
perfect.  This gives the same statement at level one.  Equivalently, one
may formulate it directly on the tame compactified moduli stack.

\begin{proposition}[Perfect Hasse-quotient pairing]
\label{prop:hasse-perfect-pairing}
For every $r$ there is a canonical perfect pairing
\[
 \langle\ ,\ \rangle_r:
 Q_r\times Q_{P+2-r}\longrightarrow \mathbb F_p.
\]
It is the product pairing followed by the trace map on the finite
Gorenstein scheme $Z$.
\end{proposition}

\begin{proof}
Because $Z$ is an effective Cartier divisor on the smooth curve $X$, its
dualizing sheaf is
\[
 \omega_Z\simeq\bigl(\omega_X\otimes\mathcal O_X(Z)\bigr)\vert_Z.
\]
Kodaira--Spencer gives $\omega_X(C)\simeq\omega^2$, while
$\mathcal O_X(Z)\simeq\omega^P$. The canonical section of
$\mathcal O_X(C)$ is nonzero on $Z$, since $C\cap Z=\varnothing$.
Consequently
\[
 \omega_Z\simeq\omega^{P+2}\vert_Z.
\]
Multiplication now gives
\[
 \mathcal Q_r\otimes\mathcal Q_{P+2-r}\longrightarrow\omega_Z.
\]
Composing global sections with
$\operatorname{Tr}_Z:H^0(Z,\omega_Z)\to\mathbb F_p$ gives the displayed
pairing. Duality for the finite Gorenstein scheme $Z$ identifies
$\mathcal Q_{P+2-r}$ with
$\mathcal Hom_Z(\mathcal Q_r,\omega_Z)$, so the pairing is perfect.
\end{proof}

Let $\Theta$ denote the Katz theta operator in characteristic $p$,
normalized by its $q$-expansion
\[
 \Theta\!\left(\sum a_nq^n\right)=\sum n a_nq^n.
\]
It raises weight by $p+1$, satisfies the product rule, and satisfies
$\Theta(A)=0$. Hence it induces
\[
 \Theta:Q_r\longrightarrow Q_{r+p+1}.
\]

\begin{proposition}[Theta adjointness]
\label{prop:theta-adjointness}
If $x\in Q_r$ and $y\in Q_{P+1-p-r}$, then
\[
 \langle\Theta x,y\rangle_{r+p+1}
   =-\langle x,\Theta y\rangle_r.
\]
More generally, for $m\geq0$ and
$y\in Q_{P+2-r-m(p+1)}$,
\[
 \langle\Theta^m x,y\rangle_{r+m(p+1)}
   =(-1)^m\langle x,\Theta^m y\rangle_r.
\]
\end{proposition}

\begin{proof}
The two products in the first formula have weight $P+2$, so both traces
are defined. The product rule gives
\[
 (\Theta x)y+x(\Theta y)=\Theta(xy).
\]
Notice that $xy$ has weight
\[
 P+1-p=(p-1)^2.
\]
Let $a$ be the canonical section on the Igusa curve, normalized by
$a^{p-1}=A$. The differential attached to $\Theta(xy)$ is
$\Theta(xy)/a^P=\Theta(xy)/A^p$ under Kodaira--Spencer, and it is exact;
this is the geometric description of theta in Gross, Proposition~5.8
(equivalently, it follows locally by differentiating the function attached
to $xy$). The residues of this differential at the supersingular points
are exactly the local components of the trace in
Proposition~\ref{prop:hasse-perfect-pairing}. An exact differential has zero
residue, so
$\operatorname{Tr}_Z(\Theta(xy))=0$. This proves the first identity.
Iteration proves the second.

For completeness, the nonreduced trace can be seen locally.  At a point
$s\in S$, choose a parameter $t$ for which $A$ is a unit times $t$.
Then $Z$ has local ring $k(s)[t]/(t^p)$, and the Grothendieck trace is the
coefficient of $t^{p-1}$, followed by
$\operatorname{Tr}_{k(s)/\mathbb F_p}$, up to a common unit fixed by the
chosen trivialization.  The coefficient of $t^{p-1}$ in the derivative of
a polynomial modulo $t^p$ is zero: it could only come from the derivative
of $t^p$, whose coefficient is $p=0$.  Thus the local trace of the exact
differential vanishes on the full thickening $pS$, not merely on its
reduced support.
\end{proof}

The weight relation relevant to the Tate complement contains a further
shift of $p(p+1)$.  Let
\[
 B:=E_{p+1}\pmod p\in H^0(X,\omega^{p+1}).
\]

\begin{lemma}[The Eisenstein bridge is a unit]
\label{lem:eisenstein-unit}
The forms $A$ and $B$ have no common zero.  Hence $B\vert_Z$ is a unit and
multiplication by $B^p$ gives an isomorphism
\[
 B^p:Q_r\xrightarrow{\ \sim\ }Q_{r+p(p+1)}
\]
for every $r$.
\end{lemma}

\begin{proof}
The Hasse invariant has a simple zero at every point of $S$.  Let
$\partial_k=\Theta-kE_2/12$ be the Serre derivative in weight $k$.  The
Ramanujan--Serre identity and the $q$-expansion principle give
\[
 12\partial_{p-1}A=B
 \qquad\text{in characteristic }p.
\]
At a zero of $A$, the connection term is a multiple of $A$ and vanishes.
Thus $\partial_{p-1}A$ is the ordinary derivative of a local equation for
$S$, up to a nonzero change of trivialization.  It is nonzero because the
zero is simple.  Therefore $B$ does not vanish on $S$.  Its restriction to
the Artinian thickening $Z$ is consequently invertible, and so is $B^p$.
\end{proof}

\begin{theorem}[Perfect Tate-complement pairing]
\label{thm:tate-complement-pairing}
For every $r$ the formula
\[
 [x,z]_r:=\langle x,B^{-p}z\rangle_r
\]
defines a canonical perfect pairing
\[
 [\ ,\ ]_r:
 Q_r\times Q_{P+2+p(p+1)-r}\longrightarrow\mathbb F_p.
\]
Under this pairing theta is skew-adjoint.  More precisely, if
\[
 z\in Q_{P+1-p+p(p+1)-r},
\]
then
\[
 [\Theta x,z]_{r+p+1}=-[x,\Theta z]_r.
\]
The same identity iterates with the sign $(-1)^m$ for $\Theta^m$.
\end{theorem}

\begin{proof}
Lemma~\ref{lem:eisenstein-unit} identifies the second factor with
$Q_{P+2-r}$, so perfectness follows from
Proposition~\ref{prop:hasse-perfect-pairing}.  Since theta is a derivation
in characteristic $p$,
\[
 \Theta(B^p)=pB^{p-1}\Theta(B)=0,
 \qquad
 \Theta(B^{-p})=0
\]
on $Z$.  Proposition~\ref{prop:theta-adjointness} therefore gives
\[
 \begin{aligned}
 {}[\Theta x,z]_{r+p+1}
 &=\langle\Theta x,B^{-p}z\rangle_{r+p+1}\\
 &=-\langle x,\Theta(B^{-p}z)\rangle_r\\
 &=-\langle x,B^{-p}\Theta z\rangle_r
 =-[x,\Theta z]_r.
 \end{aligned}
\]
Iteration proves the final assertion.
\end{proof}

\begin{lemma}[Why the undivided $p$-fold bridge fails]
\label{lem:theta-p-zero-on-hasse-quotient}
On every $Q_r$ one has $\Theta^p=0$.
\end{lemma}

\begin{proof}
Fermat's congruence and the $q$-expansion principle give the standard
identity
\[
 \Theta^p=A^{p+1}\Theta
\]
as operators on modular forms in characteristic $p$: the two sides have
the same weight and the same $q$-expansion. In the target Hasse quotient,
\[
 A^{p+1}\Theta x=A^p(A\Theta x)=0.
\]
\end{proof}

Now let
\[
 w_-=k+2i-P,
 \qquad
 w_+=k+2i^\dagger+P,
 \qquad
 i+i^\dagger=p^2+1-k.
\]
Then
\[
 w_-+w_+=2p^2+2=P+2+p(p+1),
\]
and therefore
\[
 [\ ,\ ]_{w_-}:Q_{w_-}\times Q_{w_+}\longrightarrow\mathbb F_p
\]
is a perfect pairing under which theta is skew-adjoint.  This is the
pairing at the two candidate-weight indices.  The factor $B^p$
is exactly the weight-$p(p+1)$ factor that occurs in the proved leading-block
formula for the $D=2$ obstruction.  Serre duality alone would pair weights
summing to $P+2$; the Eisenstein bridge supplies the required extra shift.
The tempting replacement by $\Theta^p$ is zero by
Lemma~\ref{lem:theta-p-zero-on-hasse-quotient}.

There is an important indexing distinction.  The obstruction to lowering
to a candidate weight $w$ naturally lies in $Q_{w+P}$, not $Q_w$; this is
proved in Proposition~\ref{prop:geometric-bockstein}.  Thus the pairing
above is a correct structural pairing, but it is not yet the pairing of the
two actual Bockstein target spaces.  The latter is constructed, with the
correct shift, in Theorem~\ref{thm:natural-bockstein-pairing}.  This
distinction prevents the constant sum of the candidate weights from being
mistaken for Bockstein compatibility.
